\chapter{Project Engineering \& Git Workflow}

\section{Software Architecture \& Directory Structure}
The repository follows a clean, modular pythonic software architecture adhering to separation of concerns:

\begin{verbatim}
dataminingproject/
+-- data/
|   +-- raw/heart_disease_risk_2026.csv
|   \-- processed/
|       +-- X_train.csv, X_test.csv, y_train.csv, y_test.csv
|       \-- model_comparison.csv
+-- src/
|   +-- utils.py                # File paths & column constants
|   +-- preprocessing.py          # ColumnTransformer & feature engineering
|   +-- statistical_tests.py      # 6 Hypothesis tests & FDR corrections
|   +-- feature_analysis.py       # Multi-perspective feature importances
|   +-- models.py                 # ML training & GridSearchCV tuning
|   +-- visualization.py          # Publication-quality figure generation
|   +-- findings.py               # Empirical discovery synthesis
|   \-- app.py                    # 9-Page Streamlit web dashboard
+-- scripts/
|   +-- download_data.py          # Automated kagglehub downloader
|   +-- train_pipeline.py         # End-to-end model training script
|   \-- generate_report_assets.py # Master figure & JSON generator
+-- models/
|   +-- best_model.joblib
|   +-- preprocessing_pipeline.joblib
|   \-- model_metadata.json
+-- results/
|   +-- statistical_tests.json
|   +-- feature_importances.json
|   \-- key_findings.json
+-- docs/
|   +-- Rapport_MiniProject.tex
|   +-- Rapport_MiniProject.pdf
|   \-- figures/
\-- tests/                        # Pytest unit tests suite
\end{verbatim}

\section{Reproducibility Protocol}
To reproduce the complete pipeline from scratch:
\begin{enumerate}
    \item Clone repository: \texttt{git clone https://github.com/hosniadilemp-a11y/DataMiningProjectExample.git}
    \item Install dependencies: \texttt{pip install -r requirements.txt}
    \item Download raw dataset: \texttt{python scripts/download\_data.py}
    \item Execute training pipeline: \texttt{python scripts/train\_pipeline.py}
    \item Generate report assets: \texttt{python scripts/generate\_report\_assets.py}
    \item Launch Streamlit app: \texttt{streamlit run src/app.py}
\end{enumerate}

\section{Git Demonstration Workflow}
\begin{tcolorbox}[colback=yellow!15!white,colframe=orange!80!black,title=\textbf{DEMONSTRATION GIT WORKFLOW NOTICE}]
This reference repository was created using \textbf{ONE GitHub account} to demonstrate a multi-student Git feature-branch workflow. The student codes \texttt{AISD05}, \texttt{SIAD19}, and \texttt{RSD20} represent \textbf{SIMULATED CONTRIBUTORS}.
\end{tcolorbox}

The repository structure incorporates four dedicated Git branches:
\begin{itemize}
    \item \texttt{main}: Production-ready stable release branch.
    \item \texttt{develop}: Integration branch for merging student feature branches.
    \item \texttt{feature/AISD05-eda}: Branch for Student AISD05 (Data Audit, EDA, Hypotheses).
    \item \texttt{feature/SIAD19-modeling}: Branch for Student SIAD19 (Preprocessing, ML Models).
    \item \texttt{feature/RSD20-streamlit}: Branch for Student RSD20 (Streamlit UI, LaTeX Report).
\end{itemize}

\section{Student Responsibility Matrix}
Table~\ref{tab:contrib_matrix} details the responsibility distribution across simulated student contributors.

\begin{table}[H]
\centering
\caption{Simulated Student Contributor Responsibility Matrix}
\label{tab:contrib_matrix}
\small
\begin{tabular}{lll}
\toprule
\textbf{Student Code} & \textbf{Primary Responsibilities} & \textbf{Git Feature Branch} \\
\midrule
AISD05 & Dataset acquisition, quality audit, EDA, hypothesis testing & \texttt{feature/AISD05-eda} \\
SIAD19 & Preprocessing pipeline, ML modeling, tuning, evaluation & \texttt{feature/SIAD19-modeling} \\
RSD20 & Streamlit web application, predictor UI, LaTeX report & \texttt{feature/RSD20-streamlit} \\
\bottomrule
\end{tabular}
\end{table}
